\begin{remark}
  \( a \) is not a simple root \( \iff (x-a)\mid f, (x-a)\mid f' \implies \gcd(f,f') \ne 1 \implies
  f\mid f'\) since \( f \) is minimal polynomial and thus irreducible. And see that \( f \mid f'
  \implies a \) is not simple root of \( f \). Now see that 
  \begin{align*}
    f&= \sum_{i=0 }^n a_i x^i \\ 
    f'&=\sum_{i=1}^n ia_i x^{i-1} 
  \end{align*}
  the degree of \( f' \) certainly goes down, then either \( f' = 0 \) or \( \deg(f') < \deg(f) \).
  Hence \( f\mid f' \iff f' = 0 \), by the fact that
  \[
    ia_i = \underbrace{i 1_K}_{=\vp(i) : \vp: \ZZ \to K} a_i
  \] 
  then see that \( f'= 0 \iff \forall\; i,  \), if \( a_i \ne 0 \implies \vp(i) = 0 \iff \text{char}(K) = p > 0 
  \) and \( f \in K[x^p] \).
\end{remark}

\begin{eg}
  Suppose that \( \text{char}(K) = p \), with \( a \in K \smallsetminus K^p \), define the polynomial
  \[
    f = x^p - a \in K[x]
  \] 
  with a root \( b \) in some algebraic closure \( \overline{K} \) of \( K \), with \( b \not\in K \). Here 
  \( b \) being a root means \( b^p = a \).
  \begin{claim}
    \( f \) is the minimal polynomial of \( b \).
  \end{claim}
  \begin{proof}
    We already know that \( f(b) = 0 \), so it is enough to show that \( f \) is \textcolor{blue}{irreducible}. 
    Consider in \( K(b)[x] \), we have 
    \[
      f = x^p - a = x^p - b^p = (x-b)^p
    \] 
    If \( f \) is not irreducible in \( K[x] \), then there exists \( g \mid f \), with \( 1 \leq \deg(g) < p \),
    where \( g \in K[x] \). Now consider putting \( g \) in \( K(b)[x] \), it is clear that to make
    \( g \) divides \( f \) there will be only one case:
    \[
      \begin{aligned}
        g & = (x-b)^i \quad 1 \leq i \leq p-1 \\ 
          &= x^i - ib x^{i-1} + \cdots 
      \end{aligned}
    \] 
    this implies that \( ib \in K \) where \( i \) is invertible in \( K \) since \( i < p \implies b \in K \),
    leading to contradiction \( \lightning \).
  \end{proof}

  \begin{conclusion}
    \( b \) is not separable over \( K \). As its minimal polynomial is \( f = x^p - a \), but \( f'
    = 0\), and thus \( b \) is not a simple root.
  \end{conclusion}
\end{eg}

\begin{eg}[\textbf{Explicit Example}]
  Let \( k \) be any field with characteristic \( p > 0 \), let \( K = k(x) \), \textbf{see that} \( x \in
  k\smallsetminus k^p \):, if \( x = (\frac{f}{g})^p \; f,g \in k[x] \), then \( xg^p = f^p \), thus
  \( p\cdot \deg(g) + 1 = p\cdot \deg(f) \), leading to contradiction to the fact that \( p \) is prime number and \( p \geq 2 \) \( \lightning \).
  \begin{note}
    Note that \( k(x) = \text{Frac}(k[x]) \).
  \end{note}
\end{eg}

\begin{definition}[\textbf{Perfect Extension}]
  A field \( K \) is \underline{perfect} if every algebraic extension \( K \hookrightarrow L \) is separable.
\end{definition}

\begin{proposition}
  A field \( K \) is perfect iff 
  \begin{itemize}
    \item either \( \text{char}(K) = 0 \).
    \item or \( \text{char}(K) = p > 0 \) and \( K = K^p \).
  \end{itemize}
\end{proposition}

\begin{proof}
  \leavevmode 
  \begin{itemize}
    \item We know in char \( 0 \), every extension is separable, thus every field of characteristic of char
      \( 0 \) is perfect. As in this case \( \gcd(f,f') = 1 \) and thus there is no multiple roots for \( f \).
    \item Suppose now \( \text{char}(K) = p \), if \( K \ne K^p \), by taking \( a\in K
      \smallsetminus K^p \) and \( b \in \overline{K} \), s.t. \( b^p a \), we saw that \( K
      \hookrightarrow K(b) \) is algebraic, but not separable, thus \( K \) is not perfect.
      Conversely, suppose now \( K = K^p \), say \( K \hookrightarrow L \) is an algebraic
      extension, suppose there exitss \( a\in L \), with minimal polynomial \( f \in K[x^p] \):
      \[
        f = a_0 + a_1 x^p + a_2 x^{2p} + \cdots + a_n x^{np}
      \]
      with \( K = K^p \implies \forall \; 0\leq i \leq n, a_i = b_i^p \) for some \( b_i \in K \),
      this implies 
      \[
        f = g^p \text{ with: } g = b_0 + b_1 x + b_2 x^2 + \cdots + b_n x^n \in K[x]
      \] 
      by using the fact that \( u \to u^p \) is ring homomorphism. Since \( p > 1 \), this leads to
      the contradiction with the irreducibility of \( f \) \( \lightning \).
  \end{itemize}
\end{proof}


\begin{eg}
  \leavevmode 
  \begin{itemize}
    \item If \( K \) is algebraically closed, then \( K \) is perfect (no non-trivial algebraic extension.)
    \item If \( K \) is a finite field, then \( K \) is perfect. 
      \begin{proof}
        If \( \text{char}(K) = p > 0 \), see 
        \begin{align*}
          F :K &\to K \\ 
          F(u) &= u^p 
        \end{align*}
        is a ring homomorphism, and it is injective, thus by finite it is also surjective. By the
        proposition above, \( K \) is perfect.
      \end{proof}
    \item If \( k \) is any field with char \( p > 0 \), then \( k(x) \ne  \) perfect.
  \end{itemize}
\end{eg}

\begin{remark}
  If \( K \hookrightarrow K_1 \hookrightarrow K_2 \) are algebraic extensions. Then \( \qo{K_2}{K} \) separable \( \implies \)
  \( \qo{K_1}{K} \) is separable (clear) and \( \qo{K_2}{K_1} \) is separable (if \( f,g \) are
  minimal polynomial of \( a \in K_2 \) over \( K, K_1 \) respectively, then \( g \mid f \) since \(
  f \in K_1[x], f(a) = 0 \implies g \) has no multi roots since \( f \) has no multi roots.)
\end{remark}

\begin{fact}
  If \( \qo{K_1}{K} \) and \( \qo{K_2}{K_1} \) are separable, then \( \qo{K_2}{K} \) is separable. \todo{Proof later.}
\end{fact}

\begin{remark}
  Also, we will need to prove that if \( L = K(a_1,\ldots, a_n) \) and \( a_1,\ldots, a_n \) are
  separable over \( K \), then \( \qo{L}{K} \) is separable.
\end{remark}

\begin{theorem}[\textbf{Primitive Element Theorem}]
  If \( K \hookrightarrow L \) is finite separable extension, then it is a \underline{simple}
  extension, i.e. \( \exists \; a \in L \), s.t. \( L = K(a) \). Here \( a \) is called a
  \underline{primitive element}.
\end{theorem}

\begin{proof}
  \leavevmode 
  \begin{itemize}
    \item If \( K \) is finite, then \( L \) is finite, by our homework problem, \( L^\times \) is cyclic,
      so we can take \( a \) to be a generator of \( L^\times \).
    \item Suppose \( K \) is infinite, and \( \qo{L}{K} \) is finite, then \( L = K(a_1,\ldots, a_n) \). We shall argue 
      by induction on \( n\geq 1 \). When \( n = 1 \) it is ok. If we know this for \( n = 2 \),
      then the inductive step is trivial, as 
      \[
        K(a_1,\ldots, a_{n-1}) = K(b) \text{ by IH } \implies L = K(b,a_n)
      \] 
      Suppose \( L = K(b,c) \), we will show: if \( \lambda \in K \) and \( a = b + \lambda c \)
      except for finitely many values of \( \lambda \), we shall get \( L = K(a) \). Now fix \(
      \lambda \in K \) and let \( a = b + \lambda c \), we have 
      \[
        K \hookrightarrow K(a) \hookrightarrow L
      \] 
      and may assume \( c\not\in K(a) \), otherwise \( b \in K(a) \implies L = K(a) \). Let \( f,g \) be the minimal polynomial 
      of \( b,c \) respectively over \( K(a) \). The roots are simple since \( \qo{L}{K} \) is separable. Let 
      \begin{align*}
        \text{roots of }f: \beta_1=b, \beta_2,\ldots, \beta_n \\ 
        \text{roots of }g: \gamma_1 = c, \gamma_2,\ldots, \gamma_n \quad (n\geq 2 \text{ since } c
        \not\in K(a)
      \end{align*}
      Let \( \overline{L} \) be some algebraic closure of \( L \), we have 
      \[
        \begin{tikzcd}
          K\arrow[hook, r] & K(a) \arrow[hook, r] \arrow[hook,rd] & K(a)(\gamma_1)
          \arrow[hook,r]\arrow[d, "\sim"] & L \arrow[dashed, "\sigma", d] \\ 
                                          & & K(a)(\gamma_2) \arrow[hook,r] & \overline{L}
        \end{tikzcd}
      \] 
      there exists morphism of extension of \( K(a) \)
      \[
        \begin{aligned}
          \sigma: L& \to \overline{L} \\ 
          \sigma(\gamma _1) &= \gamma _2
        \end{aligned}
      \] 
      \begin{note}
        \( \sigma(\beta_1 = b) = \beta_i \) for some \( i \), but note that 
        \[
          \begin{aligned}
            \sigma|_{K(a)} &= \Id  \\ 
            \implies \sigma (\underbrace{b + \lambda c}_{=\beta_i + \lambda \gamma _2=\beta _1 + \lambda \gamma _1}) &= 
            b + \lambda  c \\ 
            \implies \lambda &= \frac{\beta _i -\beta _1}{\gamma _1 - \gamma _2} \quad \text{by }
            \gamma _1 \ne \gamma _2
          \end{aligned}
        \] 
        Since \( f \) and \( g \) divide the minimal polynomial of \( b \) and \( c \) respectively
        over \( K \), then only finitely many choices for \( \beta_i, \gamma_j \implies \) only
        finitely many choices for \( \lambda \). If \( \lambda \ne  \) finitely many choices \(
        \implies  L = K(a) \).
      \end{note}
  \end{itemize}
\end{proof}

